\begin{Lemma}[McLaughlin]
\label{lem:Mac}
If $(\alpha_n, \beta_n)$ is a Bailey pair relative to $a$, then $(\alpha'_n, \beta'_n)$ is a Bailey pair relative to $a/q$, where
\begin{equation*}
\alpha'_0=\alpha_0, \qquad \alpha'_n=(1-a)\left(\frac{\alpha_n}{1-aq^{2n}}-\frac{aq^{2n-2}\alpha_{n-1}}{1-aq^{2n-2}}\right), \qquad \beta'_n=\beta_n.
\end{equation*}
\end{Lemma}

In this paper, we will show, among other things, that this lemma and the Bailey lattice can be extended to bilateral versions.


As noted in \cite{BMS} and~\cite{J}, it is possible to define for all $n\in\mathbb{Z}$ a \emph{bilateral Bailey pair} $(\alpha_n, \beta_n)$ relative to $a$ by the relation
\begin{equation}\label{bbp}
\beta_n=\sum_{j\leq n}\frac{\alpha_j}{(q)_{n-j}(aq)_{n+j}}\qquad\forall n\in\mathbb{Z}.
\end{equation}

\begin{Remark}\label{rk:bilatinversion}
The relation~\eqref{bp} defining classical (unilateral) Bailey pairs is a special instance of the above relation defining bilateral ones, as choosing $\alpha_n=0$ for negative integers $n$ in~\eqref{bbp} implies $\beta_n=0$ for $n$ negative. Actually, the converse is also true, as the classical Bailey inversion holds for bilateral Bailey pairs: $(\alpha_n, \beta_n)$ is a bilateral Bailey pair relative to $a$ if and only if
\begin{equation}\label{bilatinversion}
\alpha_n=\frac{1-aq^{2n}}{1-a}\sum_{j\leq n}\frac{(a)_{n+j}}{(q)_{n-j}}(-1)^{n-j}q^{\binom{n-j}{2}}\beta_j\qquad\forall n\in\mathbb{Z}.
\end{equation}
\textit{Thus, from all our results in this paper, one can deduce the corresponding unilateral results by setting $\alpha_n=0$ $($or equivalently $\beta_n=0)$ for all $n<0$.}
\end{Remark}

In \cite{BMS}, the Bailey lemma is extended in the following way.
\begin{Theorem}[bilateral Bailey lemma]\label{thm:bilatbaileylemma}
If $(\alpha_n, \beta_n)$ is a bilateral Bailey pair relative to $a$, then
so is $(\alpha'_n, \beta'_n)$, where
\begin{equation*}\alpha'_n=\frac{(\rho,\sigma)_n(aq/\rho\sigma)^n}{
(aq/\rho,aq/\sigma)_n}\alpha_n\qquad\mbox{and}\qquad\beta'_n=\sum_{j\leq n}\frac{(\rho,\sigma)_j(aq/\rho\sigma)_{n-j}(aq/\rho\sigma)^j}{(q)_{n-j}(aq/\rho,aq/\sigma)_n}\beta_j,
\end{equation*}
subject to convergence conditions on the sequences $\alpha_n$ and $\beta_n$, which make the relevant infinite series absolutely convergent.
\end{Theorem}

Our first result is an extension of the Bailey lattice to the bilateral case.

\begin{Theorem}[bilateral Bailey lattice]\label{thm:bilatbaileylattice}
If $(\alpha_n, \beta_n)$ is a bilateral Bailey pair relative to $a$, then~$(\alpha'_n, \beta'_n)$ is a bilateral Bailey pair relative to $a/q$, where
\begin{equation*}\alpha'_n=\frac{(\rho,\sigma)_n(a/\rho\sigma)^n}{(a/\rho,a/\sigma)_n}(1-a)\left(\frac{\alpha_n}{1-aq^{2n}}-\frac{aq^{2n-2}\alpha_{n-1}}{1-aq^{2n-2}}\right),
\end{equation*}
and
\begin{equation*}\beta'_n=\sum_{j\leq n}{(\rho,\sigma)_j(a/\rho\sigma)_{n-j}(a/\rho\sigma)^j\over (q)_{n-j}(a/\rho,a/\sigma)_n} \beta_j,
\end{equation*}
subject to convergence conditions on the sequences $\alpha_n$ and $\beta_n$, which make the relevant infinite series absolutely convergent.
\end{Theorem}

As mentioned above, several proofs have been given for the (unilateral) Bailey lattice. On the other hand, simpler proofs were given to prove the Andrews--Gordon identities without using the Bailey lattice. Here we give a very simple proof of our bilateral Bailey lattice, which when considering Bailey pairs such that $\alpha_n=0$ for $n<0$ reduces to the unilateral Bailey lattice. Hence we provide in particular a very simple proof of the classical Bailey lattice.

The key in our proof is the following simple lemma, which generalises McLaughlin's unilateral Lemma \ref{lem:Mac} and transforms bilateral Bailey pairs relative to $a$ into bilateral Bailey pairs relative to $a/q$.
\begin{Lemma}[key Lemma~1]\label{lem:key1}
If $(\alpha_n, \beta_n)$ is a bilateral Bailey pair relative to $a$, then $(\alpha'_n, \beta'_n)$ is a bilateral Bailey pair relative to $a/q$, where
\begin{equation}\label{eq:baileylattice1}
\alpha'_n=(1-a)\left(\frac{\alpha_n}{1-aq^{2n}}-\frac{aq^{2n-2}\alpha_{n-1}}{1-aq^{2n-2}}\right), \qquad \beta'_n=\beta_n,
\end{equation}
subject to convergence conditions on the sequences $\alpha_n$ and $\beta_n$, which make the relevant infinite series absolutely convergent.
\end{Lemma}
While it is not customary to do so, we give the proof in this introduction to show that it is just a few lines long and only requires the definition of bilateral Bailey pairs and elementary sum manipulations which are similar to the ones for the unilateral version.
\begin{proof}[Proof of Lemma~\ref{lem:key1}]
For all $n\in\mathbb{Z}$, we have
\begin{align*}
\sum_{j\leq n} \frac{\alpha'_j}{(q)_{n-j}(a)_{n+j}}&= \sum_{j\leq n} \frac{(1-a)}{(q)_{n-j}(a)_{n+j}}\left(\frac{\alpha_j}{1-aq^{2j}}-\frac{aq^{2j-2}\alpha_{j-1}}{1-aq^{2j-2}}\right)\\
&= \sum_{j\leq n} \frac{(1-a)\alpha_j}{(q)_{n-j}(a)_{n+j}\bigl(1-aq^{2j}\bigr)}-\sum_{j\leq n} \frac{(1-a)\bigl(1-q^{n-j}\bigr)aq^{2j}\alpha_{j}}{(q)_{n-j}(a)_{n+j+1}\bigl(1-aq^{2j}\bigr)}\\
&=\sum_{j\leq n} \frac{(1-a)\alpha_j}{(q)_{n-j}(a)_{n+j+1}\bigl(1-aq^{2j}\bigr)}\bigl(\bigl(1-aq^{n+j}\bigr)-aq^{2j}\bigl(1-q^{n-j}\bigr)\bigr)\\
&=\sum_{j\leq n} \frac{\alpha_j}{(q)_{n-j}(aq)_{n+j}}=\beta_n=\beta'_n,
\end{align*}
which is the desired result by~\eqref{bbp}.
\end{proof}
With this lemma, we can give an extremely simple proof of the bilateral Bailey lattice.
\begin{proof}[Proof of the bilateral Bailey lattice]
Start from a bilateral Bailey pair $(\alpha_n, \beta_n)$ relative to~$a$. Then apply Lemma~\ref{lem:key1} to obtain a bilateral Bailey pair $\bigl(\tilde{\alpha}_n, \tilde{\beta}_n\bigr)$ relative to $a/q$ and satisfying~\eqref{eq:baileylattice1}. Applying the bilateral Bailey lemma (see Theorem~\ref{thm:bilatbaileylemma}) to \smash{$\bigl(\tilde{\alpha}_n, \tilde{\beta}_n\bigr)$} with $a$ replaced by~$a/q$ gives the desired bilateral Bailey pair $(\alpha'_n, \beta'_n)$ relative to $a/q$.
\end{proof}

In addition to Lemma~\ref{lem:key1}, let us give a similarly simple lemma whose unilateral version is also due to McLaughlin in~\cite[Lemma~13.1\,(2)]{M} (see also Lovejoy \cite[Lemma 3.1]{Lo22}), and whose proof is very similar to the one of Lemma~\ref{lem:key1}.

\begin{Lemma}[key Lemma~2]\label{lem:key2}
If $(\alpha_n, \beta_n)$ is a bilateral Bailey pair relative to $a$, then $(\alpha'_n, \beta'_n)$ is a bilateral Bailey pair relative to $a/q$, where
\[%\label{eq:baileylattice2}
\alpha'_n=(1-a)\left(\frac{q^n\alpha_n}{1-aq^{2n}}-\frac{q^{n-1}\alpha_{n-1}}{1-aq^{2n-2}}\right), \qquad \beta'_n=q^n\beta_n,
\]
subject to convergence conditions on the sequences $\alpha_n$ and $\beta_n$, which make the relevant infinite series absolutely convergent.
\end{Lemma}

Using, as in the short proof of the bilateral Bailey lattice above, Lemma~\ref{lem:key2} followed by Theorem~\ref{thm:bilatbaileylemma} with $a$ replaced by $a/q$, we obtain the following new bilateral Bailey lattice, similar to Theorem~\ref{thm:bilatbaileylattice}. As far as we know, its unilateral version was also unknown until now.

\begin{Theorem}[new bilateral Bailey lattice]\label{thm:newbilatbaileylattice}
If $(\alpha_n, \beta_n)$ is a bilateral Bailey pair relative to~$a$, then $(\alpha'_n, \beta'_n)$ is a bilateral Bailey pair relative to $a/q$, where
\begin{equation*}\alpha'_n=\frac{(\rho,\sigma)_n(a/\rho\sigma)^n}{(a/\rho,a/\sigma)_n}(1-a)\left(\frac{q^n\alpha_n}{1-aq^{2n}}-\frac{q^{n-1}\alpha_{n-1}}{1-aq^{2n-2}}\right),
\end{equation*}
and
\begin{equation*}\beta'_n=\sum_{j\leq n}{(\rho,\sigma)_j(a/\rho\sigma)_{n-j}(a/\rho\sigma)^j\over (q)_{n-j}(a/\rho,a/\sigma)_n} q^j\beta_j,
\end{equation*}
subject to convergence conditions on the sequences $\alpha_n$ and $\beta_n$, which make the relevant infinite series absolutely convergent.
\end{Theorem}

Lemmas \ref{lem:key1} and \ref{lem:key2} can be generalised by adding an extra parameter $b$.
\begin{Lemma}[general lemma]\label{lem:baileylattices}
If $(\alpha_n, \beta_n)$ is a bilateral Bailey pair relative to $a$, then $(\alpha'_n, \beta'_n)$ is a bilateral Bailey pair relative to $a/q$, where
\[%\label{eq:baileylatticesalpha}
\alpha'_n=\frac{1-a}{1-b}\left(\frac{\bigl(1-bq^n\bigr) \alpha_n}{1-aq^{2n}}-\frac{q^{n-1}\bigl(aq^{n-1}-b\bigr) \alpha_{n-1}}{1-aq^{2n-2}}\right)
\qquad \text{and} \qquad \beta'_n=\frac{1-bq^n}{1-b}\beta_n.
\]
\end{Lemma}

\begin{Remark}
Lemma \ref{lem:key1} is the case $b=0$ and Lemma \ref{lem:key2} is the case $b \rightarrow \infty$ of Lemma \ref{lem:baileylattices}.
\end{Remark}

\begin{Remark}
While at first glance Lemma \ref{lem:baileylattices} seems more general than Lemmas \ref{lem:key1} and~\ref{lem:key2}, it is actually equivalent to these two lemmas taken together. Indeed, the bilateral Bailey pair in Lemma \ref{lem:baileylattices} is equal to $1/(1-b)$ times the bilateral Bailey pair of Lemma \ref{lem:key1} minus~${b/(1-b)}$ times the bilateral Bailey pair of Lemma \ref{lem:key2}. Using the fact that being a bilateral Bailey pair is stable under linear combination, Lemmas \ref{lem:key1} and \ref{lem:key2} imply Lemma \ref{lem:baileylattices}. Note that it is also possible to prove Lemma \ref{lem:baileylattices} directly with a similar method to the proof of Lemma \ref{lem:key1}.
\end{Remark}

Despite following from Lemmas \ref{lem:key1} and \ref{lem:key2}, this general Lemma \ref{lem:baileylattices} is still interesting as it provides in the unilateral case an ``inverse" to Lovejoy's Lemma 2.3 of \cite{Lo22}, which he first stated in \cite[equations (2.4) and (2.5)]{Lo04}. Actually, this result inspired our discovery of Lemma~\ref{lem:baileylattices}.

